% ============================================================================
%  Semantic Categorical Orchestra
%  kwasa-kwasa, formalised: the conductor of a research operating system.
%
%  Self-contained. Every object is a finite weighted graph or a construction on
%  one; only the standard external literature is cited. Interpretation is
%  confined to remarks on which no theorem depends. The paper is intended to be
%  read alone: nothing in the surrounding framework is presupposed.
% ============================================================================
\documentclass[twocolumn,11pt,a4paper]{article}

\usepackage[utf8]{inputenc}
\usepackage[T1]{fontenc}
\usepackage{lmodern}
\usepackage[a4paper,margin=2cm]{geometry}
\usepackage{amsmath,amssymb,amsthm,mathtools}
\usepackage{bm}
\usepackage{mathrsfs}
\usepackage{booktabs}
\usepackage{array}
\usepackage{enumitem}
\usepackage{microtype}
\usepackage{xcolor}
\usepackage[round,sort&compress]{natbib}
\usepackage[colorlinks=true,linkcolor=blue!45!black,
            citecolor=blue!45!black,urlcolor=blue!45!black]{hyperref}
\usepackage{cleveref}

% ---- theorem environments ----
\newtheoremstyle{sup}{6pt}{6pt}{\itshape}{}{\bfseries}{.}{.5em}{}
\theoremstyle{sup}
\newtheorem{theorem}{Theorem}[section]
\newtheorem{lemma}[theorem]{Lemma}
\newtheorem{proposition}[theorem]{Proposition}
\newtheorem{corollary}[theorem]{Corollary}
\theoremstyle{definition}
\newtheorem{definition}[theorem]{Definition}
\newtheorem{axiom}{Premise}
\newtheorem{construction}[theorem]{Construction}
\newtheorem{example}[theorem]{Example}
\theoremstyle{remark}
\newtheorem{remark}[theorem]{Remark}
\newtheorem{principle}[theorem]{Principle}

\crefname{axiom}{Premise}{Premises}
\Crefname{axiom}{Premise}{Premises}

% ---- operators and shorthand ----
\DeclareMathOperator{\Aut}{Aut}
\DeclareMathOperator*{\argmin}{arg\,min}
\DeclareMathOperator*{\argmax}{arg\,max}
\DeclareMathOperator{\hol}{hol}
\DeclareMathOperator{\dom}{dom}
\DeclareMathOperator{\cod}{cod}
\DeclareMathOperator{\im}{im}
\newcommand{\R}{\mathbb{R}}
\newcommand{\N}{\mathbb{N}}
\newcommand{\E}{E}                    % edge set
\newcommand{\reach}{\mathrm{reach}}   % reachable-state operator
\newcommand{\flo}{\beta}              % the floor
\newcommand{\med}{\mathsf{m}}         % the medium
\newcommand{\cut}{\partial}
\newcommand{\co}{\complement}
\newcommand{\str}{\sigma}             % separation cost / minimum cut
\newcommand{\Om}{\Omega}
\newcommand{\Sent}{S}                 % semantic distance functional
\newcommand{\Sflat}{S_{\flat}}        % the floor of S
\newcommand{\cell}{\mathsf{C}}        % action cell
\newcommand{\task}{t}
\newcommand{\Tasks}{T}
\newcommand{\Ens}{E}                  % ensemble of tasks
\newcommand{\intent}{\iota}           % intent
\newcommand{\mods}{\mathcal{M}}       % set of modules
\newcommand{\modu}{\mu}               % a module
\newcommand{\dsl}{\mathcal{L}}        % a module's language
\newcommand{\orch}{\mathsf{K}}        % kwasa-kwasa, the orchestrator
\newcommand{\dem}{D}                  % demand
\newcommand{\contrib}{\delta\Sent}    % contribution score
\newcommand{\Rens}{R}                 % ensemble order parameter
\newcommand{\abs}[1]{\lvert#1\rvert}
\newcommand{\norm}[1]{\lVert#1\rVert}
\newcommand{\set}[1]{\{#1\}}
\newcommand{\eps}{\varepsilon}

\title{\textbf{Semantic Categorical Orchestra}\\[3pt]
\large The Conductor of a Research Operating System:\\
A Self-Contained Formalisation of an Orchestration Language}

\author{%
Kundai Farai Sachikonye\\[2pt]
\small Technical University of Munich\\
\small Chair of Experimental Bioinformatics\\
\small \texttt{kundai.sachikonye@wzw.tum.de}}

\date{\today}

\begin{document}
\maketitle

% ============================================================================
\begin{abstract}
\noindent
We formalise an orchestration language for a research operating system. The
system presents no applications and no windows: a blank surface on which a
researcher writes a program in one language, and to which only the currently
necessary result is returned. Specialised competence---testing, simulation,
analysis, retrieval---is held not by the language but by a federation of
\emph{modules}, each a tool with its own domain-specific language that
\emph{terminates} its domain and answers within it. The orchestration language
generates no general-purpose code; it compiles intent into the modules' own
languages and routes it. This relieves the language of the burden of omniscience
and raises a different problem, which is the subject of the paper. Every module,
queried correctly, returns a \emph{correct} answer; correctness is the module's
guarantee. But a correct answer to the wrong question is worthless, and no module
can tell whether its own task was the right one to perform, because that judgment
is a property of the whole and a module is a bounded part. We prove, from a
single derived fact---a strictly positive floor on the cost of telling any part
from a non-completable whole---that necessity is unreadable from within any task
(the No-Local-Necessity theorem), that it is therefore the irreducible function
of a layer above the modules, and that this layer can certify necessity
\emph{without redoing any module's work}, by three measurements on the
task-graph alone: a \emph{contribution score} that vanishes exactly on
purposeless tasks, a \emph{holonomy} that detects when correct steps compose into
a cycle that drifts from the stated intent, and an \emph{order parameter} whose
phase transition marks the ensemble's convergence onto a common goal. We define
the language as a typed compiler from intent to module-language with a necessity
gate at its centre, prove that the gate is sound (it never prunes a necessary
task) and that orchestration converges to a region---an \emph{action cell}---of
admissible plans rather than to a single point, and we exhibit the conductor as
the one component the federation cannot do without: a hall of flawless players is
silent without it. Everything is proved of finite weighted graphs; minimum cuts
are exact, and interpretation is confined to remarks.
\end{abstract}

\noindent\textbf{Keywords:} orchestration language; domain-specific language;
research operating system; minimum cut; resolution floor; necessity;
contribution score; holonomy; order parameter; bounded observer; action cell;
separation of concerns.

\tableofcontents

% ============================================================================
\part{Foundations: the floor, parts, and the whole}
% ============================================================================

% ============================================================================
\section{Introduction}
\label{sec:intro}
% ============================================================================

\subsection{A blank surface}

Consider an operating system for research that shows nothing until asked. There
is no desktop of applications, no array of windows, no file browser; there is a
surface on which the user writes, and onto which the system returns the one piece
of information the writing requires, after which the surface is clear again. Such
a system replaces the application---a standing program waiting to be used---with
a sentence: a request, expressed once, answered once. The question this paper
addresses is what language that sentence is written in, and what that language
must \emph{be} in order for the system to work.

\subsection{Players and a conductor}

The competence of the system---the ability to test a piece of code, to simulate a
reaction, to align two sequences, to retrieve a measurement---does not live in the
language. It lives in a federation of \emph{modules}: independent tools, each
expert in one domain, each exposing its own small language by which it is
addressed. A module is to its domain what a virtuoso is to an instrument. The
orchestration language does not play the instruments; it does not generate the
low-level code that a module would otherwise need handed to it. It compiles a
request into the \emph{module's own language} and dispatches it. The competence is
the module's; the language only needs to know how to \emph{ask}.

This division has an immediate benefit and a hidden cost. The benefit: the
language need not be omniscient. It does not have to know how to write a correct
simulation, only how to ask the simulation module for one; the simulation
module's language already \emph{terminates} that domain---it admits exactly the
well-formed simulations and nothing else---so the orchestrator commits only the
intent and lets the module's compiler enforce correctness. This is the classical
separation of concerns \citep{dijkstra1972,mernik2005} carried to its conclusion.

The cost is subtler, and it is the reason the orchestrator must \emph{excel}, not
merely route. Every module, asked a well-formed question, returns a
\emph{correct} answer; that is its contract. But a correct answer to a question
that should never have been asked advances nothing. A hall of the finest players
in the world, each rendering their part flawlessly, is noise---or silence---without
someone to decide which piece is played, in what order, toward what whole. That
someone is the conductor, and the conductor's skill is not a better rendition of
any part. It is the judgment of \emph{necessity}: whether each correctly performed
task is the \emph{right} task, a \emph{needed} one, for the result the user
actually asked for. We will prove that this judgment cannot be made by any player,
and must be the defining function of the orchestrator.

\subsection{The thesis}

Three claims organise the paper, each proved below.

\emph{First} (\Cref{sec:floor,sec:nolocal}), necessity is a global property and is
invisible locally. We derive a strictly positive floor on the cost of
individuating any part of a non-completable whole, and show that a bounded module,
seeing only its own task against its own rest, cannot encode whether that task is
necessary to the whole. Correctness is local and is the module's to guarantee;
necessity is global and is not.

\emph{Second} (\Cref{sec:modules,sec:contribution,sec:holonomy,sec:order}), the
orchestrator can nonetheless certify necessity without performing any module's
work, by three measurements on the graph of tasks: a contribution score that is
zero exactly on tasks whose removal changes nothing the ensemble can achieve; a
holonomy that is nonzero exactly when correct tasks compose around a cycle that
departs from the declared intent; and an order parameter whose value, and whose
sharp transition, report how nearly the ensemble has converged on a single goal.

\emph{Third} (\Cref{sec:language,sec:cell,sec:conductor}), the orchestration
language is exactly a typed compiler from intent to module-language with a
\emph{necessity gate} at its centre. We prove the gate sound, prove that
orchestration converges to a \emph{region} of admissible plans rather than a
single plan, and prove that the conductor is the unique component the federation
cannot dispense with.

\subsection{Method}

The development is austere on purpose. Every object is a finite weighted graph or
a construction on one. ``Identity,'' ``meaning,'' ``necessity,'' ``correctness''
are each a specific graph quantity, defined before use, and minimum cuts are
computable exactly by maximum flow \citep{fordfulkerson1956,menger1927}. Where a
tempting claim could not be proved at this standard it has been removed rather
than weakened, and every interpretive reading---of a module as a player, of the
orchestrator as a conductor, of a task-cycle as a feedback loop---is confined to a
remark on which no theorem depends. The paper is self-contained: a reader who
grants only finite weighted graphs can check every claim without consulting any
other work.

% ============================================================================
\section{Contact graphs and the floor}
\label{sec:floor}
% ============================================================================

We begin with the object on which everything rests and the one constant we will
use throughout. Both are defined here from nothing.

\begin{definition}[Finite weighted graph; cut]
\label{def:graph}
A \emph{finite weighted graph} is a triple $\Gamma=(V,\E,w)$ with $V$ a finite
nonempty vertex set, $\E\subseteq\binom{V}{2}$ undirected edges, and
$w\colon\E\to\R_{>0}$. For $\varnothing\neq S\subsetneq V$ the \emph{cut} is
$\cut(S)=\set{\set{u,v}\in\E:u\in S,\ v\notin S}$, of weight
$w(\cut(S))=\sum_{e\in\cut(S)}w(e)$.
\end{definition}

\begin{definition}[Contact graph; medium; separation cost]
\label{def:contact}
A \emph{contact graph} has a distinguished vertex $\med$, the \emph{medium},
adjacent to every other vertex. The vertices $v\neq\med$ are \emph{items}. The
\emph{separation cost} of an item $v$ is the minimum weight of a cut placing $v$
opposite the medium,
\[
\str(v)=\min_{S\ni v,\ \med\notin S}w(\cut(S)),
\]
the $v$--$\med$ minimum cut, attained on finite $V$ and computable by max-flow
\citep{fordfulkerson1956,menger1927,diestel2017}. The minimising cut is the
\emph{resting cut} of $v$.
\end{definition}

\begin{remark}[The medium is ``everything else'']
\label{rem:medium}
The medium is not a place; it is the single vertex standing for all that an item
is told apart from. Its universal adjacency encodes that every item is, at
minimum, individuated against the whole rest. No theorem uses any reading of
$\med$ beyond \Cref{def:contact}.
\end{remark}

\begin{definition}[Expanding family; non-completable whole]
\label{def:noncompletable}
An \emph{expanding family} is a chain of contact graphs
$\Gamma_0\subseteq\Gamma_1\subseteq\cdots$, each adding items and contacts while
keeping the medium and all prior structure. The whole is \emph{non-completable}
if the family has no last stage: for every $\Gamma_k$ there is
$\Gamma_{k+1}\supsetneq\Gamma_k$. This is an order condition, not a cardinality
assumption.
\end{definition}

\begin{theorem}[The floor]
\label{thm:floor}
Let the whole be non-completable. Then there is a constant $\flo>0$ with
$w(e)\ge\flo$ for every contact and $\str(v)\ge\flo$ for every item: no item is
told from the rest at zero cost, and there is no cut of weight $0$ separating a
nonempty item set from its nonempty complement.
\end{theorem}

\begin{proof}
To identify an item $v$ completely is to fix it by its cut against the entire
rest of the whole. Suppose this needs zero residue: then at some finite stage
$\Gamma_k$ the whole has been brought to bear, so $\Gamma_k$ contains the whole,
contradicting non-completability (\Cref{def:noncompletable}), which supplies a
further stage with an item not yet weighed against $v$. Hence a positive residue
remains at every finite stage; set $\flo$ to its infimum. This residue is carried
by the edges of $v$'s resting cut, so each such edge has weight $\ge\flo$; every
edge lies in some singleton cut, so every weight is $\ge\flo$; and connectivity
forces $\str(v)=w(\cut(S^\ast(v)))\ge\flo>0$.
\end{proof}

\begin{remark}[Why a floor, and why it is not assumed]
\label{rem:trace}
$\flo$ is not posited; it is the trace a non-completable whole leaves on each of
its parts. A completable whole would permit zero-cost cuts and point-identities;
the impossibility of finishing the comparison is exactly what makes the floor
positive. Every later result depends only on $\flo>0$, never on its value.
\end{remark}

% ============================================================================
\section{Identity, correctness, and the action cell}
\label{sec:identity}
% ============================================================================

We now fix what ``the identity of a part'' and ``a correct result'' mean, and
prove that each is a \emph{region}, never a point. This is the structural fact
behind every later claim about necessity.

\begin{definition}[Weighted isomorphism; invariant]
\label{def:iso}
A \emph{weighted isomorphism} is a vertex bijection preserving the medium, edges,
and weights. A quantity is a \emph{graph invariant} if it agrees on isomorphic
graphs.
\end{definition}

\begin{theorem}[Identity is the invariant cut, a region]
\label{thm:region}
The separation cost and resting cut of an item are graph invariants, while its
label is permuted by automorphisms and is not. The identity of an item---the
invariant datum fixing it---is its resting cut, a set of contacts of weight
$\ge\flo>0$; in general the minimising side $S^\ast(v)$ is not the singleton
$\set{v}$. Identity is therefore a region of contacts, never a point.
\end{theorem}

\begin{proof}
An isomorphism carries each $v$--$\med$ cut to an equal-weight cut and back, so
the minimum is preserved; labels are permuted by any nontrivial automorphism.
By \Cref{thm:floor} the cut has weight $\ge\flo>0$. For non-locality, join two
dense item-clusters by a single $\flo$-weight contact: the cheapest cut against
$\med$ may peel off a whole cluster, a side of size $>1$, below the cost of
isolating any one vertex; the minimiser is then not a singleton.
\end{proof}

We carry this from a single item to a \emph{result}: the outcome of a piece of
work, scored by how far it stands from doing what was asked.

\begin{definition}[Semantic distance; the floor of a result]
\label{def:Sfunctional}
Fix a goal $g$ (a designated item, the intended terminus). The \emph{semantic
distance} of a state $x$ to $g$ is the alignment gap
\[
\Sent(x)=\frac{\str(x,g)}{\Om}\in(0,1],\qquad \Om:=w(\E),
\]
where $\str(x,g)$ is the $x$--$g$ minimum cut. By \Cref{thm:floor},
$\Sent(x)\ge\flo/\Om=:\Sflat>0$ for $x\neq g$: no state is perfectly on goal;
$\Sflat$ is the irreducible floor of $\Sent$.
\end{definition}

\begin{theorem}[Correctness is an action cell, not a point]
\label{thm:cell}
The set $\cell:=\set{x:\Sent(x)=\Sflat}$ of states at the floor---the
\emph{action cell} of the goal---is in general not a single state: distinct $x\neq
x'$ may both attain $\Sflat$. A correct result is membership in $\cell$, a
positive-volume region; it is never a unique prescribed output.
\end{theorem}

\begin{proof}
$\Sent(x)=\Sflat$ holds iff the $x$--$g$ cut attains the floor. By
\Cref{thm:region} the floor-attaining cut need not be borne by a singleton, and
two states whose cuts to $g$ have equal floor weight are $\Sent$-indistinguishable;
construct them as in \Cref{thm:region}. Hence $\abs{\cell}>1$ in general.
\end{proof}

\begin{remark}[``Correct'' is a region; this is not a defect]
\label{rem:cell}
\Cref{thm:cell} says a module that returns \emph{any} $x\in\cell$ has returned a
correct result; there is no single right output to match. This is why correctness
can be a module's local guarantee and yet leave the harder question---whether the
\emph{goal} $g$ was the one the user wanted---entirely open. We make that gap
precise next.
\end{remark}

% ============================================================================
\section{No local necessity}
\label{sec:nolocal}
% ============================================================================

This section proves the paper's pivot: a part cannot judge its own necessity. We
first separate two notions kept apart throughout.

\begin{definition}[Correctness versus necessity]
\label{def:corr-nec}
Let a \emph{task} produce a result $x$ aimed at a local goal $g_\task$. The task
is \emph{correct} if $x\in\cell(g_\task)$: it does what it was asked, to within
the floor. Given an \emph{intent} $g$ (the global goal) and an ensemble of tasks
$\Ens$, a task $\task\in\Ens$ is \emph{necessary} if removing it strictly worsens
the best semantic distance the ensemble can reach toward $g$:
\[
\min_{x\in\reach(\Ens\setminus\set\task)}\Sent_g(x)\ >\
\min_{x\in\reach(\Ens)}\Sent_g(x),
\]
where $\reach(\cdot)$ is the set of states the ensemble can produce. Correctness
is judged against $g_\task$; necessity against $g$.
\end{definition}

\begin{definition}[Bounded module; its horizon]
\label{def:bounded}
A \emph{module} is a bounded observer: a finite contact graph $\Gamma_\modu$ with
a state capacity $\kappa_\modu=\log_2\abs{V_\modu}$ bits, individuating its task
against \emph{its own} rest $\med_\modu$, with no edge to the global goal $g$
unless $g\in V_\modu$.
\end{definition}

\begin{theorem}[No local necessity]
\label{thm:nolocal}
A module cannot, in general, determine the necessity of its own task. Precisely:
if the global intent $g$ lies outside the module's graph
($g\notin V_\modu$), then necessity---a predicate on the global ensemble's reach
toward $g$---is not a function of any datum available within $\Gamma_\modu$; and
even when $g$ is present, if the ensemble's global state complexity exceeds
$\kappa_\modu$ bits, no correctness criterion encodable in the module decides it.
\end{theorem}

\begin{proof}
\emph{(i) Absent goal.} By \Cref{def:corr-nec} necessity quantifies over
$\reach(\Ens)$ and $\reach(\Ens\setminus\set\task)$, both defined by cuts toward
$g$. If $g\notin V_\modu$ there is no $g$-cut inside $\Gamma_\modu$, so no
internal quantity is a function of $\Sent_g$; two global ensembles agreeing on
$\Gamma_\modu$ but differing in whether $\task$ is necessary are indistinguishable
to the module. Hence necessity is not module-internal.

\emph{(ii) Present goal, insufficient capacity.} Suppose $g\in V_\modu$ but the
global ensemble has $N$ distinguishable necessity-configurations with $\log_2
N>\kappa_\modu$. A decision of necessity is a function from configurations to
$\set{0,1}$ that must be \emph{stored} (encoded) within the module's
$\kappa_\modu$ bits to be applied; by the pigeonhole principle a $\kappa_\modu$-bit
state cannot label $N>2^{\kappa_\modu}$ distinct configurations injectively, so
some two configurations with opposite necessity collapse to one module-state and
are decided alike---wrongly for one. No module-encodable criterion is correct on
all of them.
\end{proof}

\begin{corollary}[Necessity is the layer above]
\label{cor:above}
Necessity is decidable only by an observer whose graph contains the global goal
$g$ and the whole ensemble $\Ens$, with capacity $\ge\log_2 N$. Since a module is
by definition bounded to its own task, this observer is a distinct layer
\emph{above} the modules.
\end{corollary}

\begin{proof}
Immediate from \Cref{thm:nolocal}: the two failure modes are exactly absence of
$g$ and insufficiency of capacity; an observer holding $g$, $\Ens$, and enough
bits suffers neither, and is not any single module.
\end{proof}

\begin{remark}[The orchestra, stated exactly]
\label{rem:orchestra}
\Cref{thm:nolocal} is the formal content of the orchestra. Each player renders a
correct part (membership in its own action cell, \Cref{thm:cell}); no player can
hear whether the \emph{piece} is right, because the piece is a property of the
whole ensemble's convergence on a goal that no part contains
(\Cref{cor:above}). The conductor is not a better player; the conductor is the
observer the theorem requires. No theorem depends on the reading.
\end{remark}

% ============================================================================
\part{The federation and the three measurements}
% ============================================================================

% ============================================================================
\section{Modules, languages, and the orchestrator}
\label{sec:modules}
% ============================================================================

We now build the federation and the layer above it, and fix what each is allowed
to do.

\begin{definition}[Module; its language; termination of a domain]
\label{def:module}
A \emph{module} $\modu$ is a pair $(\dsl_\modu,\;[\![\cdot]\!]_\modu)$ of a
\emph{domain-specific language} $\dsl_\modu$---a set of well-formed
\emph{instructions}---and an interpretation taking each instruction to a result.
The module \emph{terminates} its domain: $\dsl_\modu$ admits exactly the
well-formed tasks of the domain and no others, and every admitted instruction is
interpreted to a correct result ($x\in\cell(g_\task)$, \Cref{thm:cell}). A module
is a bounded observer (\Cref{def:bounded}).
\end{definition}

\begin{definition}[Intent; the orchestrator]
\label{def:orchestrator}
An \emph{intent} $\intent$ is a request written by the user in a single
\emph{orchestration language}, naming a global goal $g$. The \emph{orchestrator}
$\orch$ (the conductor; \emph{kwasa-kwasa}) is a map that takes an intent and a
federation $\mods=\set{\modu_1,\dots,\modu_k}$ and returns a result surfaced to
the user. $\orch$ generates \emph{no} domain instructions of its own beyond
emitting, for each chosen subtask, a well-formed instruction in the relevant
module's language $\dsl_\modu$. The competence of each domain is the module's;
$\orch$'s competence is elsewhere, and is the subject of \Cref{sec:contribution}
onward.
\end{definition}

\begin{principle}[Separation of competence]
\label{prin:separation}
Two competences are kept disjoint. \emph{Correctness} is each module's, enforced
by $\dsl_\modu$'s termination of its domain (\Cref{def:module}). \emph{Necessity}
is the orchestrator's, and by \Cref{thm:nolocal} can be nowhere else. The
orchestrator never re-checks a module's correctness, and a module never judges its
own necessity.
\end{principle}

\begin{remark}[Why the orchestrator emits a language, not code]
\label{rem:emit-dsl}
A general code generator must hold the competence of every target domain---it must
know how to write a correct simulation, a correct test, a correct alignment---which
makes it an all-knowing part, contradicting boundedness and inviting error at
every domain it touches. Emitting into $\dsl_\modu$ removes this: the orchestrator
commits only the intent (which task, toward which goal), and the module's
termination of its domain supplies correctness. The orchestrator is thereby freed
to spend its whole capacity on the one thing no module can do---necessity.
This is the classical compiler discipline of lowering to a typed target
\citep{aho2006,mernik2005} put to the service of \Cref{prin:separation}; no
theorem depends on the reading.
\end{remark}

\begin{definition}[Task graph]
\label{def:taskgraph}
For an intent $\intent$ the orchestrator forms a \emph{task graph}: a directed
graph $G_\Ens=(\Ens,A)$ whose vertices are tasks (each routed to a module) and
whose arcs $u\to v$ record that $v$ consumes $u$'s result. Sources are the
intent's givens; the designated sink carries the global goal $g$. Each task $\task$
has a local goal $g_\task$ and, once run, a result $x_\task\in\cell(g_\task)$.
\end{definition}

The next three sections give the orchestrator's three instruments, each a
measurement on $G_\Ens$ that needs no module re-execution.

% ============================================================================
\section{Contribution: the necessity of a single task}
\label{sec:contribution}
% ============================================================================

\begin{definition}[Ensemble floor; contribution score]
\label{def:contribution}
The \emph{ensemble floor} is the best semantic distance the ensemble can reach
toward $g$,
\[
\Sflat(\Ens)=\min_{x\in\reach(\Ens)}\Sent_g(x)\ \ge\ \Sflat .
\]
The \emph{contribution} of a task $\task$ is the increase in the ensemble floor
caused by removing it,
\[
\contrib(\task,\Ens)=\Sflat(\Ens\setminus\set\task)-\Sflat(\Ens)\ \ge 0 .
\]
\end{definition}

\begin{theorem}[Contribution decides necessity]
\label{thm:contribution}
$\contrib(\task,\Ens)>0$ iff $\task$ is necessary
(\Cref{def:corr-nec}); $\contrib(\task,\Ens)=0$ iff $\task$ is
\emph{purposeless}---its removal leaves unchanged everything the ensemble can
achieve toward $g$. Monotonicity holds: removing a task cannot lower the floor,
so $\contrib\ge0$ always.
\end{theorem}

\begin{proof}
Removing a task can only shrink $\reach$ (fewer producible states), so the minimum
over a subset cannot fall: $\Sflat(\Ens\setminus\set\task)\ge\Sflat(\Ens)$, giving
$\contrib\ge0$. By \Cref{def:corr-nec} necessity is exactly
$\Sflat(\Ens\setminus\set\task)>\Sflat(\Ens)$, i.e.\ $\contrib>0$; its negation
$\contrib=0$ is equality of the floors, i.e.\ purposelessness.
\end{proof}

\begin{corollary}[A correct task may be purposeless]
\label{cor:correct-purposeless}
There exist tasks with $x_\task\in\cell(g_\task)$ (perfectly correct) and
$\contrib(\task,\Ens)=0$ (purposeless): a flawless answer that the ensemble did
not need. Correctness does not imply necessity.
\end{corollary}

\begin{proof}
Let $\task$ produce a correct result on a local goal $g_\task$ disconnected from
the sink $g$ in $G_\Ens$ (its arc does not reach the sink). Then $\reach$ toward
$g$ is unchanged by its presence, so $\contrib=0$ while $x_\task\in\cell(g_\task)$.
\end{proof}

\begin{remark}[The pruning the conductor performs]
\label{rem:prune}
\Cref{thm:contribution,cor:correct-purposeless} give the conductor its first act:
compute $\contrib$ for each task and \emph{prune} those with $\contrib=0$. A
purposeless task is correct and removable; pruning it changes nothing the ensemble
can achieve and spares its cost. This is necessity exercised positively, by
subtraction. Soundness---that pruning never removes a necessary task---is
\Cref{thm:gate-sound}.
\end{remark}

% ============================================================================
\section{Holonomy: necessity around a cycle}
\label{sec:holonomy}
% ============================================================================

Single-task contribution does not catch a subtler failure: tasks each correct and
each necessary, composed around a feedback cycle that drifts from the goal. We
measure the drift directly.

\begin{definition}[Transport; specification; holonomy]
\label{def:holonomy}
Each arc $u\to v$ of $G_\Ens$ carries a \emph{transport} $T_{u\to v}$: the map by
which $v$ transforms $u$'s result. For a directed cycle $c=(v_0\to
v_1\to\cdots\to v_n=v_0)$ the \emph{accumulated transport} is
$T_c=T_{v_{n-1}\to v_0}\circ\cdots\circ T_{v_0\to v_1}$. The intent supplies a
\emph{specification} $T_c^{\mathrm{spec}}$, the transformation the cycle is
declared to realise (the identity if the cycle should restore its starting
state). The \emph{holonomy} of $c$ at a probe $x$ is
\[
\hol(c,x)=\norm{T_c(x)-T_c^{\mathrm{spec}}(x)} .
\]
\end{definition}

\begin{theorem}[Cycle inconsistency]
\label{thm:holonomy}
If $\hol(c,x)>0$ for some reachable probe $x$ on some cycle $c$ of $G_\Ens$, then
the ensemble is globally inconsistent with the intent: there is no assignment of
results, each locally correct, whose composition around $c$ matches the
declared specification. Zero holonomy on every cycle is necessary for global
consistency; on a DAG (no cycles) holonomy is vacuous and consistency is decided
by the contribution scores alone.
\end{theorem}

\begin{proof}
Composition around $c$ returns to $v_0$; consistency with the intent requires the
returned state equal the specified one, i.e.\ $T_c(x)=T_c^{\mathrm{spec}}(x)$ for
all reachable $x$, i.e.\ $\hol(c,x)=0$. A single $x$ with $\hol(c,x)>0$ exhibits a
state the cycle transforms against specification; since each $T_{u\to v}$ is fixed
by its (correct) module, no reassignment of \emph{correct} results removes the
mismatch---it lives in the composition, not in any node. Necessity of the
condition is its definition; on a DAG there is no cycle, so the universally
quantified condition is vacuously true and only \Cref{thm:contribution} remains.
\end{proof}

\begin{remark}[The route audit]
\label{rem:route}
Holonomy is an audit of the \emph{route}, not the endpoints. Every node may be
correct and every node necessary, yet the loop they form may accumulate a drift
that no node reveals---a discrepancy visible only by transporting a probe all the
way around and comparing with what was promised. Kirchhoff's loop law
\citep{kirchhoff1847} and the holonomy of a connection
\citep{kobayashinomizu1963} are the same shape: a quantity that is consistent iff
its accumulation around every closed loop vanishes. The conductor's second act is
to enumerate the cycles (by Johnson's algorithm \citep{johnson1975}, after a
strongly-connected decomposition \citep{tarjan1972}) and test each for holonomy.
No theorem depends on the reading.
\end{remark}

% ============================================================================
\section{The order parameter: convergence on one goal}
\label{sec:order}
% ============================================================================

The third instrument measures, in a single scalar, how nearly the whole ensemble
is pulling toward one goal, and detects the sharp moment at which it locks on.

\begin{definition}[Phase; coupling; order parameter]
\label{def:order}
Assign each task $\task$ a \emph{phase} $\theta_\task\in[0,2\pi)$: the direction,
in goal-space, of its result's pull (its alignment angle toward $g$). Arcs couple
phases with strength $K$ proportional to shared contribution. The \emph{ensemble
order parameter} is
\[
\Rens\,e^{\mathrm i\psi}=\frac1{\abs\Ens}\sum_{\task\in\Ens}e^{\mathrm
i\theta_\task},\qquad \Rens\in[0,1],
\]
the magnitude of the mean phase: $\Rens=0$ when pulls cancel (no shared goal),
$\Rens=1$ when all pulls align (one goal) \citep{kuramoto1975,kuramoto1984,
strogatz2000}.
\end{definition}

\begin{theorem}[Regimes and the locking transition]
\label{thm:order}
Under phase coupling on $G_\Ens$ there is a critical coupling $K_c$ such that for
$K<K_c$ the only stable state is incoherent ($\Rens\approx0$: the tasks do not
share a goal), and for $K>K_c$ a coherent state with $\Rens>0$ emerges and grows.
The order parameter stratifies the ensemble into regimes from \emph{turbulent}
($\Rens$ small, no common action cell) through \emph{coherent} ($\Rens$ near $1$,
converging on a common cell) to \emph{phase-locked} ($\Rens=1$); the onset at
$K_c$ is a genuine transition in the sense of \citet{landau1937}.
\end{theorem}

\begin{proof}[Proof sketch]
The mean-field reduction of phase coupling \citep{kuramoto1984,strogatz2000} gives
a self-consistency equation $\Rens=K\,\Rens\,F(\Rens)$ with $F$ the
coupling-weighted alignment kernel; $\Rens=0$ always solves it, and a second,
positive solution bifurcates from it exactly when $K$ exceeds
$K_c=1/F(0^+)$. Below $K_c$ only $\Rens=0$ is stable; above, the coherent branch
is stable and increasing in $K$. The free-energy expansion in the order parameter
near $K_c$ is that of a second-order transition \citep{landau1937}, giving the
sharp onset. Full coherence $\Rens=1$ is the fixed point at which all phases
coincide.
\end{proof}

\begin{remark}[A map, not a verdict]
\label{rem:regimemap}
The conductor's third act returns not a pass/fail but a \emph{regime map}:
$(\Rens,\Sflat(\Ens),\hol,\contrib)$. By \Cref{thm:cell} ``correct'' is already a
region, so the honest report of an ensemble's health is likewise a region-valued
characterisation---how coherently the tasks converge---rather than a point verdict.
A low $\Rens$ says the ensemble is doing many correct things that do not add up to
one goal; that is precisely a failure of \emph{necessity}, made quantitative. No
theorem depends on the reading.
\end{remark}

% ============================================================================
\part{The language, the cell of plans, and the conductor}
% ============================================================================

% ============================================================================
\section{kwasa-kwasa as a typed compiler with a necessity gate}
\label{sec:language}
% ============================================================================

We now assemble the three instruments into the orchestration language and prove
its central property: it judges necessity soundly while never touching
correctness.

\begin{definition}[The orchestration pipeline]
\label{def:pipeline}
$\orch$ maps an intent $\intent$ (global goal $g$) to a surfaced result by:
\begin{enumerate}[label=\textup{(\arabic*)},leftmargin=2em,itemsep=1pt]
\item \emph{Decompose}: form a task graph $G_\Ens$ (\Cref{def:taskgraph}),
routing each task to a module whose domain it falls in.
\item \emph{Lower}: emit, for each task, a well-formed instruction in its module's
language $\dsl_\modu$ (\Cref{rem:emit-dsl}); correctness is the module's
(\Cref{def:module}).
\item \emph{Gate}: judge necessity on $G_\Ens$---prune tasks with $\contrib=0$
(\Cref{thm:contribution}), reject ensembles with $\hol>0$ on any cycle
(\Cref{thm:holonomy}), and report the regime $\Rens$ (\Cref{thm:order}).
\item \emph{Surface}: return the result of the gated, coherent ensemble; clear the
surface.
\end{enumerate}
\end{definition}

\begin{definition}[Well-typed orchestration]
\label{def:welltyped}
A task is \emph{well-typed} if its instruction is well-formed in its module's
language and its arcs connect output types to input types along $G_\Ens$. $\orch$
\emph{accepts} an ensemble iff it is well-typed, has zero holonomy on every cycle,
and contains no purposeless task after pruning.
\end{definition}

\begin{theorem}[Soundness of the necessity gate]
\label{thm:gate-sound}
The gate is \emph{sound}: it never prunes a necessary task and never rejects a
globally consistent ensemble. Formally, if $\task$ is necessary
(\Cref{def:corr-nec}) then the gate retains it; if $\Ens$ is globally consistent
with the intent (zero holonomy on every cycle, every task necessary) then the gate
accepts it.
\end{theorem}

\begin{proof}
The gate prunes $\task$ only when $\contrib(\task,\Ens)=0$, which by
\Cref{thm:contribution} is exactly purposelessness, the negation of necessity; so
a necessary task ($\contrib>0$) is never pruned. The gate rejects only on
$\hol>0$ for some cycle, which by \Cref{thm:holonomy} is exactly global
inconsistency; so a consistent ensemble (all holonomies zero) is never rejected.
Both gate actions are characterisations, not heuristics, so each is sound by the
theorem it invokes.
\end{proof}

\begin{theorem}[The gate touches only necessity, never correctness]
\label{thm:gate-orthogonal}
The gate's three measurements---$\contrib$, $\hol$, $\Rens$---are functions of the
task graph $G_\Ens$ and the results $x_\task$ alone, and never re-evaluate whether
$x_\task\in\cell(g_\task)$. Correctness is read from the module's contract and
treated as given; the gate decides only necessity and consistency.
\end{theorem}

\begin{proof}
$\contrib$ depends on $\reach$ and $\Sflat(\Ens)$ (\Cref{def:contribution}); $\hol$ on
the transports and the specification (\Cref{def:holonomy}); $\Rens$ on the phases
(\Cref{def:order}). Each is computed from $G_\Ens$ and the delivered results; none
recomputes a module's local goal $g_\task$ or re-checks $x_\task\in\cell(g_\task)$,
which is supplied by \Cref{def:module}. Hence the gate is orthogonal to
correctness, realising \Cref{prin:separation}.
\end{proof}

\begin{remark}[Why the language can be thin]
\label{rem:thin}
\Cref{thm:gate-orthogonal} is the reason the orchestration language is small. It
carries none of the domains' competence---no module's grammar, no domain
correctness---only the three graph measurements and the routing. A user learns to
\emph{express intent and route it}, not to write simulations or tests; the surface
of the language is the entire interface to the system, and it is thin precisely
because competence lives below it. No theorem depends on the reading.
\end{remark}

% ============================================================================
\section{Orchestration converges to a region of plans}
\label{sec:cell}
% ============================================================================

We show that the gate does not seek a unique plan---there is none to seek---but
relaxes the ensemble into an \emph{action cell of plans}, and that this relaxation
halts or declines.

\begin{definition}[Plan; demand; relaxation]
\label{def:relax}
A \emph{plan} is an accepted ensemble (\Cref{def:welltyped}). Its \emph{demand} is
$\dem(\Ens)=\Sflat(\Ens)-\Sflat$, the ensemble's distance above the irreducible
floor. The \emph{relaxation} repeatedly: prunes a purposeless task, repairs a
holonomy violation by re-routing the offending cycle, or adds a task that strictly
lowers $\Sflat(\Ens)$---each step a committing change to $G_\Ens$.
\end{definition}

\begin{theorem}[Monotone relaxation; quiescence or decline]
\label{thm:relax}
Along the relaxation the demand $\dem(\Ens)$ is non-increasing and bounded below by
$0$. The relaxation either reaches \emph{quiescence}---a plan whose demand cannot
be lowered, the ensemble resting in the action cell of the intent
(\Cref{thm:cell})---in finitely many steps, or it does not halt, in which case no
plan attains the intent and the honest output is to \emph{decline} rather than
surface a non-converged ensemble.
\end{theorem}

\begin{proof}
Each step either removes a task that does not raise $\Sflat$ (pruning,
$\dem$ unchanged or lower), repairs a cycle (restoring consistency without raising
$\Sflat$), or adds a task that strictly lowers $\Sflat$; none raises $\dem$, so
$\dem$ is non-increasing, and $\dem\ge0$ by \Cref{def:Sfunctional}. The space of
ensembles over finitely many available tasks and routings is finite, so a
non-increasing demand either attains its infimum at a finite step---quiescence, a
fixed point in the sense of \citet{banach1922,tarski1955,knastertarski1928}---or
the relaxation cycles without lowering it, never reaching the floor. In the latter
case $\dem(\Ens)>0$ for every reachable plan, so no plan lands in $\cell$;
surfacing any of them would report attainment that does not hold, and declining is
the only output consistent with the floor.
\end{proof}

\begin{theorem}[The plan is a region, certified blind]
\label{thm:plan-region}
At quiescence the set of admissible plans is in general a positive-volume region,
not a single plan: distinct accepted ensembles may share the minimal demand. The
orchestrator certifies a plan by the \emph{absence of any lowering move}---a
property of the demand---without ever reading the meaning of any module's result.
\end{theorem}

\begin{proof}
By \Cref{thm:cell} the floor of $\Sent$ is attained on a region; distinct ensembles
reaching that region have equal minimal demand and are gate-indistinguishable, so
the admissible set has $>1$ element in general. Acceptance tests only
$\contrib,\hol,\Rens$ (\Cref{thm:gate-orthogonal}), all functions of the graph and
the delivered results, never of an interpreted ``meaning'' of a result; hence the
certification is blind.
\end{proof}

\begin{remark}[Why blind certification is the right standard]
\label{rem:blind}
The orchestrator cannot read what a module's result \emph{means}---that would
require occupying the module's entire domain, i.e.\ being the module, collapsing
\Cref{prin:separation}. It does not need to. Necessity and consistency are
graph-level facts about contribution, holonomy, and coherence, decidable without
interpretation. The conductor judges whether the parts \emph{add up}, not what each
part privately is. No theorem depends on the reading.
\end{remark}

% ============================================================================
\section{The conductor is indispensable}
\label{sec:conductor}
% ============================================================================

We close with the structural statement that motivates the whole construction: the
federation cannot do without the orchestrator, and the orchestrator's
indispensable function is exactly necessity.

\begin{theorem}[Indispensability of the conductor]
\label{thm:indispensable}
Let a federation $\mods$ of modules be given, each individually flawless (every
delivered result correct, \Cref{def:module}). Without the orchestration layer:
\textup{(i)} purposeless tasks are undetectable, since no module can compute
$\contrib$ over the global ensemble (\Cref{thm:nolocal}); \textup{(ii)} cyclic
inconsistency is undetectable, since holonomy lives in the composition and in no
node (\Cref{thm:holonomy}); \textup{(iii)} convergence on a single goal is
unmeasured, since $\Rens$ is a property of the ensemble, not of any task
(\Cref{def:order}). Hence a federation of correct modules, however excellent,
cannot certify that its work is necessary, consistent, or convergent: it can be a
hall of flawless players producing no coherent result.
\end{theorem}

\begin{proof}
Each of (i)--(iii) is a global quantity---a function of $G_\Ens$ and $g$---and by
\Cref{thm:nolocal,cor:above} no bounded module computes a global necessity
predicate. A federation is exactly a set of bounded modules; absent a layer
holding $g$ and $\Ens$, none of the three measurements is available, so none of
the three failures is detectable. Correctness of every part is consistent with all
three failures by \Cref{cor:correct-purposeless} and \Cref{thm:holonomy}. Thus the
federation alone cannot certify its own coherence.
\end{proof}

\begin{theorem}[Master equivalence]
\label{thm:master}
For a federation orchestrated toward a non-completable whole, the following are
equivalent:
\begin{enumerate}[label=\textup{(\roman*)},leftmargin=2em,itemsep=1pt]
\item there is a floor $\flo>0$ on individuating any part of the whole
(\Cref{thm:floor});
\item correctness is a region (an action cell), not a point
(\Cref{thm:cell});
\item necessity is not decidable within any bounded module
(\Cref{thm:nolocal});
\item judging necessity, consistency, and convergence is the irreducible function
of a layer above the modules---the conductor (\Cref{cor:above,thm:indispensable});
\item the orchestration language is a thin, typed compiler that lowers intent into
module languages and gates only on the three graph measurements, certifying plans
blind and to within the floor
(\Cref{thm:gate-sound,thm:gate-orthogonal,thm:plan-region}).
\end{enumerate}
Each is the same fact: telling a part from a non-completable whole costs a
positive, irreducible residue; that residue makes correctness a region and
necessity global; a global judgment cannot be made by a part; so it is made by the
conductor, whose language need carry nothing but the judgment.
\end{theorem}

\begin{proof}
$(\mathrm i)\!\Rightarrow\!(\mathrm{ii})$ is \Cref{thm:cell};
$(\mathrm{ii})\!\Rightarrow\!(\mathrm{iii})$ holds because a region-valued
correctness leaves the global goal underdetermined locally, the content of
\Cref{thm:nolocal}; $(\mathrm{iii})\!\Rightarrow\!(\mathrm{iv})$ is
\Cref{cor:above} with \Cref{thm:indispensable};
$(\mathrm{iv})\!\Rightarrow\!(\mathrm v)$ is the construction of
\Cref{sec:language,sec:cell}; and $(\mathrm v)\!\Rightarrow\!(\mathrm i)$ because
the gate's measurements are defined through $\Sflat=\flo/\Om>0$, whose positivity
is the floor. The cycle closes.
\end{proof}

\subsection{Computational note}

Every quantity above is computable on finite weighted graphs without
interpretation: minimum cuts (hence $\str$, $\Sent$, $\Sflat$, $\contrib$) by
maximum flow \citep{fordfulkerson1956,menger1927}; cycles for holonomy by
strongly-connected decomposition \citep{tarjan1972} and elementary-circuit
enumeration \citep{johnson1975}; the order parameter by the mean-phase sum
(\Cref{def:order}); and the regime by the mean-field threshold
(\Cref{thm:order}). The contribution score requires, in the worst case, one
ensemble-floor evaluation per task (ablation), which the static measurements
$\hol$ and $\Rens$ do not; a regime map can thus be produced before paying the
ablation cost.

\subsection{Scope}

Every result is a theorem about finite weighted graphs and the constructions on
them, provable from \Cref{def:graph} and the derived floor \Cref{thm:floor}. The
development establishes \emph{that} correctness is local and region-valued,
\emph{that} necessity is global and not locally decidable, \emph{that} an
above-the-modules conductor is required and suffices, and \emph{that} the
orchestration language is exactly a thin gated compiler. It does not prescribe the
internal design of any module or the syntax of any module language; those are the
federation's, by \Cref{prin:separation}. One hypothesis beyond the graphs is used:
that the whole is non-completable (\Cref{def:noncompletable}), an order condition,
not a cardinality posit, and exactly what forces the floor; absent it, the floor
may vanish and the separation of competence is not compelled. Interpretive
readings---of modules as players, of the orchestrator as a conductor, of a blank
surface as the user interface, of holonomy as a feedback drift---are confined to
remarks on which no theorem depends, and material that could not be proved at this
standard has been omitted.

\subsection{Conclusion}

A research operating system that shows nothing until asked needs one language to
be asked in, and that language need not be wise. Wisdom---the competence of each
domain---belongs to a federation of modules, each terminating its domain in its own
tongue and answering correctly within it. What such a federation cannot supply,
from a theorem rather than an oversight, is the judgment of whether its correct
answers were the \emph{needed} ones: necessity is a property of the whole, and a
bounded part cannot read the whole. The orchestration language is therefore not a
better player but the conductor: it lowers a researcher's intent into the modules'
own languages, and then, on the graph of tasks alone, prunes what contributes
nothing, rejects what drifts around a cycle, and reports how nearly the work
converges on one goal---certifying the plan blind, to within the floor, never to a
point. Its excellence is concentrated where no module can help: in telling the
right work from the merely correct. A hall of the finest players is silent without
it; with it, the silence of the surface between requests is the sign that exactly
the necessary work, and no more, was done.

% ============================================================================
\bibliographystyle{plainnat}
\bibliography{references}
% ============================================================================

\end{document}